\section{Post-Quantum Security}
\label{sec:pq}
%% ============================================================

The Lux ecosystem has achieved comprehensive post-quantum (PQ)
coverage across all cryptographic layers. This section documents the
deployed PQ stack and the Cloud HSM integration that protects master
key material.

\subsection{PQ Scorecard}

\begin{table}[H]
\centering
\caption{Complete post-quantum scorecard. All layers are PQ-safe
  except EVM chain signing (immutable protocol constraint).}
\label{tab:pq-scorecard}
\begin{tabular}{@{}llll@{}}
\toprule
\textbf{Layer} & \textbf{Algorithm} & \textbf{PQ Status} & \textbf{Status} \\
\midrule
Disk encryption     & AES-256 (sqlcipher)          & Safe (128-bit PQ)    & Deployed \\
HKDF derivation     & HKDF-SHA-256                 & Safe (128-bit PQ)    & Deployed \\
Field encryption    & AES-256-GCM                  & Safe (128-bit PQ)    & Deployed \\
S3 backup           & age ML-KEM-768+X25519 (\texttt{age1pq}) & Safe (FIPS 203) & Deployed \\
TLS                 & X25519MLKEM768 (first curve) & Safe (hybrid PQ)     & Deployed \\
JWT signing         & ML-DSA-65 (FIPS 204)         & Safe (Module-LWE+SIS)& Deployed \\
JWKS validation     & ML-DSA-65                    & Safe                 & Deployed \\
MPC transport       & Hybrid KEM (X25519+ML-KEM-768) & Safe              & Deployed \\
MPC key shares      & PQ KEM encryption            & Safe                 & Deployed \\
Master keys         & Cloud HSM (FIPS 140-2 L3)    & Hardware isolation   & Deployed \\
Chain signing       & secp256k1 ECDSA              & \textbf{Not PQ-safe} & EVM constraint \\
\bottomrule
\end{tabular}
\end{table}

\noindent The only remaining non-PQ component is EVM chain signing
(secp256k1 ECDSA), which is an immutable constraint of the Ethereum
Virtual Machine specification. All other layers are PQ-safe.

\subsection{S3 Backup Encryption (Deployed)}

All age encryption now uses ML-KEM-768 + X25519 hybrid recipients,
available natively in age v1.3.0+ via the \texttt{HybridRecipient}
construction. Recipient public keys use the \texttt{age1pq} prefix.

Key encapsulation produces two shared secrets $ss_1$ (X25519) and
$ss_2$ (ML-KEM-768), combined via:
\begin{equation}
  \text{file\_key} = \text{HKDF-SHA256}(ss_1 \| ss_2, \texttt{"age-hybrid-v1"})
\end{equation}
An attacker must break \emph{both} X25519 and ML-KEM-768 to recover
the file key. Label enforcement prevents mixing PQ and classical
recipients (no downgrade).

\subsection{TLS Post-Quantum (Deployed)}

All TLS 1.3 connections use X25519MLKEM768 as the first curve in the
supported groups list. Servers and clients negotiate the hybrid
key exchange, providing PQ protection for data in transit. Classical
X25519 remains as fallback for clients that do not support PQ TLS.

\subsection{JWT Signing with ML-DSA-65 (Deployed)}

Hanzo IAM issues JWT tokens signed with ML-DSA-65 (FIPS~204),
replacing ECDSA. The JWKS endpoint exposes the ML-DSA-65 public key.
All services validate JWTs using the PQ-safe ML-DSA-65 signature.

ML-DSA-65 provides EUF-CMA security under the Module-LWE and
Module-SIS hardness assumptions, ensuring that a quantum adversary
cannot forge valid JWT tokens.

\subsection{MPC Transport and Key Shares (Deployed)}

MPC node-to-node communication uses hybrid KEM (X25519 + ML-KEM-768)
for key agreement. MPC key shares at rest are encrypted with PQ KEM
encryption, ensuring that captured key shares remain confidential
against future quantum attacks.

\subsection{Cloud HSM for Master Keys (Deployed)}

Master encryption keys are stored in Cloud HSM (GCP Cloud KMS,
FIPS~140-2 Level~3 certified Cavium HSMs). The master key never
leaves the HSM boundary:

\begin{itemize}[nosep]
  \item \textbf{Wrap}: Application sends plaintext DEK to HSM API;
    HSM encrypts with master key internally; returns wrapped DEK.
  \item \textbf{Unwrap}: Application sends wrapped DEK to HSM API;
    HSM decrypts with master key internally; returns plaintext DEK.
  \item \textbf{Sign}: Application sends message hash to HSM API;
    HSM signs with ML-DSA-65 private key internally; returns signature.
\end{itemize}

Three keyrings per ecosystem, 12 keys total:
\begin{itemize}[nosep]
  \item \texttt{lux-devnet}, \texttt{lux-testnet}, \texttt{lux-mainnet}
  \item Mainnet keys use HSM protection level (FIPS 140-2 L3);
    devnet/testnet use SOFTWARE protection for cost efficiency.
\end{itemize}

Key rotation creates a new key version. Old wrapped DEKs remain
decryptable with the old version. New wrap operations use the latest
version. Formal proofs of HSM key isolation, wrap/unwrap correctness,
and key rotation are in \texttt{proofs/lean/Crypto/HSM.lean}.

\subsection{Harvest-Now-Decrypt-Later: Closed}

The combination of ML-KEM-768 hybrid encryption (S3 backups),
X25519MLKEM768 TLS (transit), and ML-DSA-65 JWT signing (auth)
closes the harvest-now-decrypt-later attack vector. An adversary
who captures encrypted data today cannot decrypt it with a future
quantum computer, because every layer uses a PQ-safe algorithm.

%% ============================================================
